\chapter{总结与展望}

\section{主要工作}
本文围绕地下停车场应急定位问题，完成了从模型构建到算法验证的完整流程，主要工作包括：
\begin{enumerate}
  \item 建立 BLE RSSI 观测与步行约束融合模型；
  \item 提出鲁棒加权与滑动窗口联合优化方法；
  \item 设计初始化、异常处理和可行域投影机制；
  \item 在多噪声、多遮挡条件下完成实验评估。
\end{enumerate}

\section{主要结论}
实验结论可总结为：
\begin{itemize}
  \item 融合模型显著优于单一 RSSI 定位；
  \item 运动约束可有效抑制轨迹抖动和跳点；
  \item 鲁棒权重机制在遮挡场景中表现稳定；
  \item 在课程级数据规模下可满足实时计算需求。
\end{itemize}

\section{不足}
当前工作仍有不足：尚未在真实地下停车场长周期实测数据上验证；对人群密集和多目标互扰场景考虑不充分；尚未引入多层立体停车场复杂拓扑。

\section{后续工作}
后续可从以下方向推进：
\begin{itemize}
  \item 采集真实应急演练数据进行实证评估；
  \item 融合 IMU/UWB 形成多源因子图定位；
  \item 建立多目标协同定位与调度接口；
  \item 开发可视化指挥端并完成闭环联调。
\end{itemize}
